% Avsnitt 15.8, ur lectures/L15/appendix/b_exercises.md.

\section{Övningar}
\label{c15:sec:exercises}

\begin{aside}
De här övningarna befäster kapitlet. Konstruera din egen \file{rx_shift_reg.vhd} utifrån
specifikationen i \secref{c15:sec:iface} till \ref{c15:sec:trace}; den utdelade testbänken, körd
enligt bilaga~\ref{app:sim}, är kontrollen av att den beter sig korrekt.
\end{aside}

\exercise{\code{rx_shift_reg}}{VHDL}
\label{c15:ex:rxsr}
Skriv din egen \file{rx_shift_reg.vhd} utifrån kapitlet. Verifiera den:

\begin{shell}
cd controller
ghdl -a --std=93 can_def.vhd rx_shift_reg.vhd rx_shift_reg_tb.vhd
ghdl -e --std=93 rx_shift_reg_tb
ghdl -r --std=93 rx_shift_reg_tb --assert-level=error
\end{shell}

Om en kontroll fallerar, läs assertion-meddelandet; det namnger exakt vilket sampel det gäller och
vad som förväntades.

\exercise{Ett medvetet stoppbitsbrott}{Simulering}
\label{c15:ex:violation}
\xpart{a} Konstruera (på papper) en 6-bitars sekvens på ledningen som utgör ett äkta brott mot
stoppbitsregeln: fem lika riktiga bitar i rad, följda av en sjätte lika bit där en stoppbit skulle ha
brutit följden.

\xpart{b} Mata in din sekvens från \textbf{a)} i \code{rx_shift_reg} (antingen genom att utöka
\file{rx_shift_reg_tb.vhd} eller genom att skriva en kort ny testbänk) och bekräfta att
\code{stuff_error} pulsar på exakt den sjätte biten, inte tidigare.

\xpart{c} \code{rx_shift_reg} försöker inte återhämta sig eller omsynkronisera efter ett
\code{stuff_error}; den rapporterar det bara under en cykel och fortsätter att ackumulera som om
ingenting hänt. Förklara, med en eller två meningar, varför det är en godtagbar design här, givet hur
\code{can_controller} (\chapref{c17:ch}) kommer att reagera på det.

\exercise{Högerjusteringen, och det som ligger ovanför den}{Simulering}
\label{c15:ex:justify}
\Secref{c15:sec:iface} säger att en färdig grupp hamnar \textbf{högerjusterad i sina lägsta
\code{bit_count} bitar}, och att \code{can_controller} förlitar sig på det när det skivar upp
rekonstruerade fält. Test 3 och 4 i den utdelade testbänken prövar det: en 6-bitars grupp, sedan en
3-bitars grupp, utan reset emellan, med \code{bit_count} ändrat mellan grupperna på precis det sätt
\code{can_controller} gör mellan ett fält och nästa.

Anta en nyss nollställd \code{rx_shift_reg}, \code{enable = '1'} och inga stoppbitar någonstans
(ingen följd når någonsin fem). Processen i \secref{c15:sec:process} skiftar in varje accepterad bit
i registrets minst signifikanta ände och nollställer aldrig registret mellan grupper; anta att din
gör detsamma. Den utdelade testbänken kräver egentligen inte det, och del \textbf{d)} återkommer
till varför.

\xpart{a} Mata med \code{bit_count = 6} in de sex bitarna \code{0 1 0 1 0 1} (i den ordningen, en per
\code{sample}). Skriv ut alla åtta bitarna i \code{data} när \code{valid} pulsar. Vad står i
\code{data(7 downto 6)}, och varför?

\xpart{b} Sätt \code{bit_count = 3} utan att nollställa, och mata in \code{0 1 1}. Skriv ut alla åtta
bitarna i \code{data} när \code{valid} pulsar igen. Vad står nu i \code{data(7 downto 3)}?

\xpart{c} Ditt svar på \textbf{b)} visar att bitarna ovanför \code{bit_count} håller rester från den
\emph{föregående} gruppen. Förklara varför det är ofarligt här, och vad \code{can_controller} måste
göra (och aldrig får göra) när det läser \code{rxsr_data} för ett fält som är smalare än 8 bitar.

\xpart{d} \code{rx_shift_reg} håller två ganska olika sorters tillstånd: skiftregistret som
ackumulerar den aktuella gruppen, och följdräknaren (\code{consecutive}, \code{last_bit}) som avgör
var stoppbitar hör hemma. Att nollställa dem är inte lika säkert i båda fallen.
\begin{itemize}
  \item Anta att modulen nollställde sitt \textbf{skiftregister} vid varje \code{valid}. Vilket av
    dina svar ovan skulle ändras, och skulle något fält någonsin bli \emph{felläst} som en följd?
    Notera att den utdelade testbänken bara kontrollerar de lägsta \code{bit_count} bitarna i
    \code{data}, så båda designerna passerar den.
  \item Anta att den nollställde \textbf{följdräknaren} i samma ögonblick. Namnge vad som går sönder,
    och säg vilken regel från \chapref{c14:ch} det speglar.
\end{itemize}

Skillnaden mellan de två svaren är skälet till att \secref{c15:sub:defaults} säger att följdräknaren
nollställs bara av reset och inte säger någonting alls om skiftregistret.
